\begin{table}[H]
\centering
\caption{Productividad laboral agregada a partir del PIB y la GEIH, 2010--2025}
\label{tab:pib_geih_productividad_total}
\small
\begin{tabular}{lrrr}
\toprule
Indicador & 2010 & 2025 & Crec. anualizado \\
\midrule
PIB por trabajador (Millones de pesos de 2015 por ocupado) & 38,0 & 44,5 & 1,05\% \\
PIB por hora trabajada (Pesos de 2015 por hora) & 15.625 & 19.563 & 1,51\% \\
\bottomrule
\end{tabular}
\caption*{\footnotesize Nota: el PIB se expresa en pesos constantes de 2015 y se obtiene como la suma de los cuatro trimestres de cada año. El PIB por trabajador se calcula dividiendo el PIB anual por los ocupados expandidos de la GEIH. El PIB por hora divide el PIB anual por las horas anuales trabajadas, estimadas como horas semanales ponderadas por 52. Fuente: cálculos propios con DANE, PIB trimestral por el enfoque de producción, y GEIH.}
\end{table}